\documentclass[german,
               fontsize=12pt,
               oneside,
               BCOR=0mm, 
               DIV=9,
               bibliography=totoc,
               numbers=noendperiod,
               pagesize,
               normalheadings,
               headsepline,
               parskip=full,
               ]{scrartcl}

% fonts & Symbols
\usepackage[T1]{fontenc}
\usepackage{lmodern, pifont, marvosym}
\renewcommand*\familydefault{\sfdefault}
\usepackage[scaled=.85]{beramono}
%\renewcommand{\ttdefault}{fvm}

% line spacing
\usepackage{setspace}

% miscellaneous spacing and offset measures
\setlength{\fboxsep}{0mm}
\setlength{\parindent}{0mm}
\setlength{\belowcaptionskip}{0em} 

% color
\usepackage{color,layout,float}
\definecolor{htwk}{RGB}{35,71,160}
\definecolor{graphicbg}{RGB}{240,240,240}
\definecolor{listingbg}{RGB}{240,240,240}
\definecolor{rulegray}{RGB}{125,125,125}
\definecolor{todo}{RGB}{220,0,0}
\definecolor{comment}{RGB}{63,127,95}
\definecolor{darkgray}{RGB}{180,180,180}

% encoding and language
\usepackage[utf8]{inputenc}
\usepackage[ngerman]{babel}

% theorems
\usepackage[framed]{ntheorem}
\usepackage{framed}
\theorempreskipamount-0.3em
\theorempostskipamount-0.2em
\renewcommand*\FrameCommand{\fboxsep10pt \fcolorbox{listingbg}{listingbg}}
%\renewcommand*\FrameCommand{{\color{black}\vrule width 2pt \hspace{1em}}}
\theoremstyle{nonumberplain}
\theoremsymbol{}
\theoremseparator{:}
\newframedtheorem{definition}{Definition}

% heading & footer
\usepackage{scrpage2}
\automark{section}
\pagestyle{scrheadings}
\setheadwidth[0pt]{\linewidth}

% figures
\usepackage[pdftex]{graphicx}
\graphicspath{{figures/}{images/}{gfx/}{../figures/}}

% include entire pdf pages
\usepackage[final]{pdfpages}

% TOC, LOT, LOF, LOL layout style
\usepackage{tocstyle}
\usetocstyle{standard}

% Title format
\usepackage{titlesec}
\titleformat{\paragraph}[display]{\bfseries\normalsize}{}{}{}[]
\titlespacing*{\paragraph}{0pt}{0pt}{-1\parskip}[0pt]

% code listings
\usepackage{listings}
\lstset{captionpos           =b,
        basicstyle           =\ttfamily \footnotesize,
        keywordstyle         =\bfseries,
        aboveskip            =\parskip,
        breaklines           =true,
        breakautoindent      =false,
        tabsize              =4,
        numbers              =none, 
        frame                =single,
        framerule            =0pt,
        framesep             =.5em,
        xleftmargin          =.5em,
        framexleftmargin     =0em,
        backgroundcolor      =\color{listingbg},
        showstringspaces     =false}%
\renewcommand{\lstlistlistingname}{Listingverzeichnis}

% itemize, enumerate, description
\usepackage{enumitem}
\setlist[1]{leftmargin=*}
\setlist{labelsep=.8em}
\setitemize{label=\Forward}
	
% captions
\usepackage[small, bf]{caption}
\usepackage{amsmath}
\usepackage{amsfonts}
\usepackage{amssymb}

\AtBeginDocument{\numberwithin{lstlisting}{section}}
\AtBeginDocument{\numberwithin{figure}{section}}

% phonetic spelling
\usepackage[T1]{tipa}

% comment environment
\usepackage{verbatim}

% hyperlinks and meta-information
\usepackage[bookmarksnumbered,
            linktoc             =all,
            colorlinks          =false,
            urlcolor            =htwk,
            linkcolor           =htwk,
            citecolor           =htwk,
            pdfborder           ={0 0 0},
            linkbordercolor     ={1 1 1},
            urlbordercolor      ={1 1 1},
            pdfstartview        ={Fit},
            pdfpagelayout       ={SinglePage},
            pdfauthor           ={\AUTHOR},
            pdftitle            ={\TITLE},
            pdfkeywords         ={\KEYWORDS},
            pdfsubject          ={\SUBJECT},
            ]{hyperref}

% footnote numbering
%\usepackage{chngcntr} 
%\counterwithout{footnote}{chapter}
            
% list of abbreviations, glossary
\usepackage[
   nonumberlist,   
   acronym,        
   toc,            
   section=section 
  ]{glossaries}  
\makeglossaries
%\input{acronym.tex}
%\input{glossary.tex}

\newglossarystyle{customlistdotted}{%
\glossarystyle{altlist}          % base this style on the listdotted style
\renewcommand{\glsgroupskip}{}      % make nothing happen between groups
}

% page numbering 
\pagenumbering{arabic}

\AtBeginDocument{\author{\AUTHOR}}
\AtBeginDocument{\title{\TITLE}}
\AtBeginDocument{\subtitle{\SUBTITLE}}

%%%% bibliography %%%%%%%%%%%%%%%%%%%%%%%%%%%%%%%%%%%%%%%%%%%%%%%%%%%%%%%%%%%%%%%%%%%%%%%%%% 

\usepackage[babel,german=quotes]{csquotes}
\usepackage[style=alphabetic,
            hyperref=true,
            ]{biblatex}
%  colon after author name
\renewcommand*{\labelnamepunct}{\addcolon\space}

% heading bibliography
\DefineBibliographyStrings{ngerman}{
    %bibliography={Literaturverzeichnis}
    references={Literaturverzeichnis}    % document class: scrartcl
}
% acronyms (simulate small caps for latin modern font)
\renewcommand*{\mkbibacro}[1]{
    {\textsf{\footnotesize{#1}}}}

% name format: "Lastname, F."
\DeclareNameFormat{default}{%
  \iffirstinits
    {\usebibmacro{name:last-first}{#1}{#4}{#5}{#7}}
    {\usebibmacro{name:last-first}{#1}{#4}{#5}{#7}}%
  \usebibmacro{name:andothers}}

% "/" between multiple author names 
\renewcommand*{\multinamedelim}{\addslash}
\renewcommand*{\finalnamedelim}{\addslash}

% change abbreviation urlseen for web-references
\DefineBibliographyStrings {german}{%											
urlseen = {Letzter Zugriff:},
}

\bibliography{bibliography}
%%%% custom commands %%%%%%%%%%%%%%%%%%%%%%%%%%%%%%%%%%%%%%%%%%%%%%%%%%%%%%%%%%%%%%%%%%%%%% 

\newcommand{\todo}[1]{{\bfseries \color{todo}TODO: \color{comment}#1}}
\newcommand{\emptypage}{\newpage\thispagestyle{empty}\mbox{}\newpage}

% titlepage header
\usepackage[absolute, overlay]{textpos}
\newcommand{\logowidth}{20mm}
\newcommand{\marginwidth}{.5\paperwidth-.5\linewidth}
\newcommand{\logoxpos}{(\marginwidth+.5\linewidth-\logowidth/2)+\logowidth/20}
\newcommand{\titlepageheader}{\begin{textblock*}{\linewidth}(\marginwidth,35mm)
                                 {\color{htwk}
                                 \rule{\linewidth}{1.5mm}}
                              \end{textblock*}
                              \begin{textblock*}{\paperwidth}(\logoxpos,22mm)
                                 \begin{figure}[H]
                                    \includegraphics[width=\logowidth]{htwk_logo}
                                 \end{figure}
                              \end{textblock*}%
}

% formatting and naming conventions
\usepackage{xspace}
\newcommand{\triplets}{\textit{Triplets}\xspace}
\newcommand{\uniform}{\textit{Uniform}\xspace}
\newcommand{\nextfit}{\textit{Next-Fit}\xspace}
\newcommand{\firstfit}{\textit{First-Fit}\xspace}
\newcommand{\bestfit}{\textit{Best-Fit}\xspace}
